% !TeX root = ../thesis.tex

%%%%%%%%%%%%%%%%%%%%%%%%%%%%%%%%%%%%%%%%%%
\thesischapter[ch:discussion][0.6\textwidth]
  {``On the whole, scientific methods are at least as important results of investigation as any other results.''}
  {--- Friedrich Nietzsche, 1878, translated from German}
  {Discussion and conclusion}
%``Im Ganzen sind die wissenschaftlichen Methoden mindestens ein ebenso wichtiges Ergebniss der Forschung als irgend ein sonstiges Resultat''
\nocite{nietzsche1878human}
\chapterabstract{This chapter synthesizes the two main contributions of the thesis, \texttt{NetCoMi} for microbial association network analysis and \texttt{permApprox} for improved permutation-based inference. It discusses how both developments emerged from practical challenges in constructing and comparing microbial networks and how they relate to a broader microbiome analysis workflow. Using the PASTURE application as an illustration, the chapter reflects on the complementary roles of community-level, feature-level, and network-based analyses, and on the importance of preprocessing choices, methodological transparency, and sensitivity analyses. Finally, it discusses limitations of the current software implementations, possible future extensions of both packages, and the challenges of interpreting results from sparse and compositional microbiome data.}
\chaptercontents
%%%%%%%%%%%%%%%%%%%%%%%%%%%%%%%%%%%%%%%%%%

%%%%%%%%%%%%%%%%%%%%%%%%%%%%%%%%%%%%%
\section{Bringing the contributions together}
%%%%%%%%%%%%%%%%%%%%%%%%%%%%%%%%%%%%%

This thesis addressed statistical challenges in microbiome data analysis at several levels, ranging from community summaries and feature-level comparisons to microbial association networks. Its two main methodological contributions may initially appear only loosely connected. \texttt{NetCoMi} provides an integrated framework for constructing, analyzing, visualizing, and comparing microbial association networks, whereas \texttt{permApprox} improves the resolution of permutation p-values when only a limited number of permutations is computationally feasible. Their connection emerged directly from the practical requirements of comparative network analysis: testing differences in associations or network properties requires repeated reconstruction of networks and can therefore become computationally expensive, while the resulting edge-wise and node-wise tests often entail substantial multiplicity burdens.

The development of \texttt{NetCoMi} was motivated by the observation that microbial network analysis is rarely based on a single method. Instead, it requires a chain of decisions concerning data filtering, zero handling, normalization or transformation, association estimation, sparsification, the representation of edges, network characterization, visualization, and, where relevant, network comparison. The package integrates these steps into a coherent workflow while retaining methodological flexibility. Its contribution is therefore not to prescribe one universally correct network construction procedure, but to make established methods accessible within a common framework and to expose the choices that define the resulting network.

The limitations of empirical permutation p-values became apparent during the development and application of comparative network procedures in \texttt{NetCoMi}. In particular, the number of permutations that can be generated is constrained when each permutation requires repeated association estimation or network construction. The resulting discrete p-value resolution can be insufficient for multiple-testing-adjusted inference. \texttt{permApprox} addresses this problem by providing a general framework for tail approximation that avoids zero p-values caused by bounded fitted GPD tails. Although its initial motivation originated from network comparison, the application in this thesis demonstrates that the same issue arises throughout a broader microbiome analysis workflow, including community-level and feature-level analyses.


%%%%%%%%%%%%%%%%%%%%%%%%%%%%%%%%%%%%%
\section{An integrated view of microbiome analysis}
%%%%%%%%%%%%%%%%%%%%%%%%%%%%%%%%%%%%%

The PASTURE application was deliberately not restricted to microbial networks. Instead, it illustrated an analysis workflow that began with descriptive characterization and community-level summaries, continued with feature-level analyses, and subsequently considered association networks as an additional perspective on the data. This ordering is important. Network analysis should not be treated as an isolated endpoint that replaces other microbiome analyses. Rather, it can complement diversity analysis, ordination, differential abundance analysis, and differential distribution analysis by addressing questions about the structure of associations between taxa.

The application also illustrates why these analysis targets should not be expected to yield identical findings. Genera that differed in abundance or distribution between EBF groups were not necessarily identified as central taxa or hubs in the group-specific networks. Conversely, taxa with notable structural positions in the networks did not necessarily stand out in the feature-level analyses. This is not a contradiction. Differential abundance and differential distribution analyses concern marginal taxon-level behavior, whereas network analysis concerns patterns of co-variation or conditional dependence among taxa. Similarly, a differential association between two taxa does not imply that either taxon is differentially abundant, and differential abundance of two taxa does not imply that their association differs between groups.

These distinctions have practical consequences for interpretation. Evidence that converges across several analysis targets may strengthen a biological hypothesis, but discordant findings can also be informative. For example, a taxon may change in abundance without changing its position in the estimated association structure, or an association may differ between groups without either taxon exhibiting a pronounced marginal shift. An integrated workflow should therefore not search for one privileged result, but should use the different analyses to characterize complementary aspects of microbial community variation.

%%%%%%%%%%%%%%%%%%%%%%%%%%%%%%%%%%%%%
\section{Lessons for microbiome analysis}
%%%%%%%%%%%%%%%%%%%%%%%%%%%%%%%%%%%%%

A first lesson is that preprocessing choices are part of the statistical analysis rather than preliminary technical details. Taxonomic aggregation, filtering, zero replacement, normalization, and transformation affect the data representation on which all subsequent analyses operate. The PASTURE application used a largely consistent preprocessing strategy across several analyses to support comparability. Nevertheless, not all methods can use exactly the same representation, because their assumptions and input requirements differ. This became particularly apparent in the network comparison, where the differential association analysis and the group-specific network construction required method choices tailored to the respective inferential targets. Consistency is therefore desirable where possible, but it should not be pursued at the expense of using a method outside its intended setting.

A second lesson is that flexibility requires transparency. The availability of many plausible methods is scientifically valuable because it enables sensitivity analyses and adaptation to the characteristics of a specific data set. At the same time, it creates substantial researcher degrees of freedom. As demonstrated by \citet{ullmann2023overoptimism}, trying several plausible workflows in unsupervised microbiome analysis and selectively reporting the result that appears most favorable can produce over-optimistic conclusions. This risk applies directly to network learning, where alternative choices for filtering, association estimation, sparsification, and network representation can lead to different apparent hubs, modules, and group differences.

Consequently, microbial network analyses should document the complete construction workflow, including the treatment of zeros, the association measure, the sparsification procedure, the selected network type, and relevant visualization choices. When several workflows are scientifically plausible, sensitivity analyses should be preferred over post hoc selection of one preferred network. In confirmatory settings, important analytical decisions should ideally be specified before inspecting the final network results. Software frameworks such as \texttt{NetCoMi} can facilitate this transparency, but they cannot replace methodological judgment.

A third lesson concerns permutation-based inference. Permutation tests are attractive in microbiome analysis because they can be applied to complex statistics for which tractable parametric null distributions are unavailable. However, the validity of a permutation procedure depends on the exchangeability assumptions implied by the null hypothesis and on repeating those parts of the analysis workflow that are affected by group allocation. Moreover, a limited number of permutations imposes a finite resolution on empirical p-values. This limitation becomes particularly consequential in high-dimensional analyses with multiple testing, where small p-values are required to retain discoveries after adjustment. The results of this thesis show that tail approximation can be useful in this setting, but it should be regarded as a refinement of, rather than a replacement for, carefully designed permutation procedures.

%%%%%%%%%%%%%%%%%%%%%%%%%%%%%%%%%%%%%
\section{Limitations and future directions}
%%%%%%%%%%%%%%%%%%%%%%%%%%%%%%%%%%%%%

The contributions of this thesis provide practical advances for microbial network analysis and permutation-based inference, but they also leave several methodological and software-related questions open. Some of these limitations arise from the scope of the current implementations, whereas others reflect more fundamental challenges of sparse and compositional microbiome data. The following discussion first considers possible extensions of \texttt{NetCoMi} and \texttt{permApprox}, before returning to limitations of the PASTURE application.

For \texttt{NetCoMi}, an important limitation concerns the functionality currently provided by the package. Although the package was designed to cover a wide range of methods for the different steps of the network analysis workflow, this collection cannot be complete, particularly in a rapidly evolving methodological field. We have addressed this limitation by allowing users to provide an externally computed association matrix and thereby apply the subsequent workflow steps also to association estimators that are not implemented in \texttt{NetCoMi}. This flexibility is useful, but it currently has an important consequence for comparative inference: permutation-based testing is not available when the association matrix is supplied externally, because the association estimation step cannot be repeated within the permutation procedure. Developing a general interface that allows externally specified association-estimation procedures to be incorporated into permutation-based network comparison would therefore be an important extension.

A second limitation is that \texttt{NetCoMi} currently focuses on networks within a single microbial domain. In particular, multi-domain networks involving, for example, bacteria and fungi are not directly supported throughout the workflow. The \texttt{SpiecEasi} package, for example, already supports the estimation of cross-domain association matrices when data from several microbial domains are provided. Such matrices can currently be supplied to \texttt{NetCoMi} externally, but the remaining workflow, especially network visualization, is not specifically designed for multi-domain settings and may require substantial manual customization. Future development could provide dedicated support for domain-specific node attributes, cross-domain edge representations, and visualization strategies that make the structure of multi-domain networks easier to inspect and interpret.

A further limitation concerns comparative analyses. The inferential functionality currently focuses on the comparison of two group-specific networks. This covers many common microbiome applications, but it does not directly address settings with more than two groups, repeated measurements, or longitudinal data. Several networks can already be inspected descriptively, for example by reusing a common layout to facilitate visual comparison across time points or study groups. However, corresponding formal procedures for comparing multiple networks or testing changes in network structure over time are not currently available. Extending the framework to repeated-measures and longitudinal designs would therefore be particularly valuable.

Beyond comparisons between observed groups, another promising extension would be the comparison of an estimated microbial network with ensembles of reference networks. For example, one could generate a large number of Erd\H{o}s--R\'enyi or Barab\'asi--Albert networks with suitable size and density characteristics and assess whether observed global or local network properties differ from those expected under these reference models. Such comparisons should be based on distributions of reference networks rather than on a single simulated realization. In the application, the Graphlet Correlation Matrices of the observed networks were compared visually with those obtained from such standard network models. A simulation-based reference approach would make it possible to move beyond this descriptive comparison and formally assess whether the observed graphlet correlation structure differs from that expected under a specified generative model. The resulting reference distributions could furthermore be used to test other network properties, including connectivity, edge density, average path length, or modularity. Overall, they could provide a more principled way to characterize whether an observed microbial network exhibits structural features that are unusual relative to a specified null or generative model.

These software-related limitations are distinct from the more fundamental limitations of network interpretation. Estimated edges represent statistical associations or conditional dependencies and should not be interpreted automatically as direct biological interactions, causal effects, or ecological mechanisms. They may reflect shared environmental conditions, host-related factors, latent confounding, taxonomic aggregation, or residual artifacts induced by sequencing and preprocessing. Likewise, centrality measures and hub nodes describe properties of an estimated network, not confirmed biological importance or causal influence.

For \texttt{permApprox}, the central limitation is the reliance on an adequate approximation of the relevant permutation-distribution tail. The current GPD-based approach assumes that the selected tail can be described sufficiently well by a fitted GPD above an appropriate threshold. When the available tail is highly irregular, contains too few exceedances, or fails the corresponding goodness-of-fit test, empirical permutation $p$-values remain the fallback. This is deliberately conservative and may provide limited resolution in precisely those settings where small $p$-values are required after multiple-testing adjustment. Future work could therefore investigate additional approximation families and fallback strategies for situations in which the GPD fit is not sufficiently reliable.

Moreover, the meaningful use of \texttt{permApprox} currently requires users to assess whether the selected settings are appropriate for the permutation distribution at hand. This includes, for example, the side of the test, the choice of null centering, and the compatibility of the fitted tail with the shape and modality of the permutation distribution. Although this flexibility is useful for adapting the method to different applications, more automated diagnostic checks could help identify settings that are unlikely to yield meaningful approximations. Dedicated plotting functions for \texttt{permApprox} objects could further support this assessment by providing standardized visualizations of the permutation distribution, threshold selection, fitted tail, and resulting empirical and approximated $p$-values.

Finally, the PASTURE application illustrates several challenges that arise when applying a broad microbiome analysis workflow to sparse sequencing data. Even after taxonomic aggregation and prevalence filtering, many genera were characterized by substantial proportions of zero counts. This affects not only association estimation, but also the reliability and interpretation of transformations, feature-level procedures, and diversity-related summaries. The different analyses in the application were selected to address complementary aspects of the data, but their results remain conditional on the chosen preprocessing strategy, the retained taxa, and the assumptions of the respective methods.

The application should therefore be interpreted primarily as an illustration of how community-level, feature-level, and network-based analyses can be combined within one coherent workflow. It does not aim to provide a comprehensive epidemiological analysis of the PASTURE cohort or definitive evidence for stable ecological mechanisms underlying the observed EBF-related patterns. Although the cohort contains substantially richer information, only a deliberately restricted subset of samples and variables was considered here in order to demonstrate the proposed workflow in a focused and reproducible setting. Future analyses could build on the full data set to investigate additional exposures, covariates, and longitudinal aspects of microbial community development.

%%%%%%%%%%%%%%%%%%%%%%%%%%%%%%%%%%%%%
\section{Conclusion}
%%%%%%%%%%%%%%%%%%%%%%%%%%%%%%%%%%%%%

This thesis began with a practical problem in microbial network analysis: how can the many methodological steps required to construct, analyze, and compare microbial association networks be combined into a coherent and reproducible workflow? This question motivated the development of \texttt{NetCoMi}, but it also led to a broader perspective on microbiome analysis. Rather than treating preprocessing, estimation, visualization, and inference as separate technical tasks, this thesis considers them as connected parts of one analytical workflow.

The PASTURE application illustrated how community-level summaries, feature-level analyses, and network-based approaches can complement one another without necessarily identifying the same taxa or yielding identical conclusions. The two main contributions of this thesis address different parts of this workflow: \texttt{NetCoMi} provides an integrated framework for microbial association network analysis, whereas \texttt{permApprox} improves permutation-based inference when computational constraints limit the number of feasible permutations. Together, they link practical software implementation with statistical methodology.

Reliable microbiome analysis does not result from applying the most elaborate method, producing the most attractive network plot, or obtaining the smallest possible $p$-value. It requires transparent analytical choices, sensitivity analyses where appropriate, and cautious interpretation in light of compositionality, sparsity, and limited information in sequencing data. The methods developed in this thesis are intended to support this process by making complex analyses more coherent, reproducible, and statistically defensible.
